\section{ERS/SRS final}
\label{sec:ers-final}

\subsection{Alcance de esta secci\'on}

El ERS/SRS es un entregable independiente de este informe: reside completo en \texttt{01\_ERS/} del repositorio, en su versi\'on \textbf{4.0} del 22 de agosto de 2026, etiquetada como \texttt{v4.0-PE5} en el historial del repositorio (\S6.1), y se compila por separado. Esta secci\'on no lo reproduce. Presenta su estructura conforme a ISO/IEC/IEEE 29148:2018 (Tabla~\ref{tab:estructura-ers}), el cat\'alogo de requisitos clave, y lo que cambi\'o en la versi\'on 4.0 respecto de la 3.0 entregada en la Unidad III.

\subsection{Estructura conforme a ISO/IEC/IEEE 29148:2018}

\begin{table}[H]
\centering
\small
\begin{tabular}{|l|p{10.2cm}|}
\hline
\cellcolor{sigagreen}\textcolor{white}{\textbf{Secci\'on del ERS}} & \cellcolor{sigagreen}\textcolor{white}{\textbf{Contenido}} \\ \hline
I. Introducci\'on & Prop\'osito, alcance del producto, glosario de 23 t\'erminos, referencias y visi\'on general \\ \hline
II. Descripci\'on general & Perspectiva y l\'imite del sistema, funciones, 5 clases de usuario, mapa de 8 \emph{stakeholders} con matriz poder/inter\'es, 8 conflictos potenciales, entorno operativo, 7 suposiciones, 9 dependencias y modelado i* (SD y SR) \\ \hline
III. Requisitos espec\'ificos & Interfaces externas con 4 mockups, 25 requisitos funcionales con plantilla de 8 atributos, 17 requisitos no funcionales cuantificados, requisitos legales LOPDP, 18 restricciones de dise\~no y 17 historias de usuario con criterios Gherkin \\ \hline
IV. Modelado UML & Casos de uso general, especificaci\'on textual de 17 casos de uso, clases refinadas, secuencia y actividad por caso de uso, estados, componentes, despliegue y DFD de nivel 1 \\ \hline
V. Priorizaci\'on y trazabilidad & Clasificaci\'on de requisitos, MoSCoW justificado, modelo de Kano, WSJF y matriz de trazabilidad extendida \\ \hline
VI. Producto M\'inimo Viable & Alcance planificado, desviaci\'on, arquitectura, despliegue y estado verificado contra c\'odigo \\ \hline
VII. Componente emp\'irico & Protocolo, PICOC, hip\'otesis, plan de an\'alisis estad\'istico y registro \\ \hline
VIII--X & Conclusiones, referencias y anexos de plan, evidencias y declaraci\'on de uso de IA \\ \hline
\end{tabular}
\caption{Estructura del ERS/SRS versi\'on 4.0}
\label{tab:estructura-ers}
\end{table}

\subsection{Cat\'alogo de requisitos funcionales}

Veinticinco requisitos funcionales, priorizados en 17 Must y 8 Should (Tabla~\ref{tab:catalogo-rf}). Cada uno se especifica en el ERS con ocho atributos: identificador, descripci\'on, actor y origen, entradas y salidas, precondici\'on y postcondici\'on, prioridad MoSCoW y criterio de verificaci\'on cuantificado.

\begin{table}[H]
\centering
\small
\begin{tabular}{|l|p{9.3cm}|l|}
\hline
\cellcolor{sigagreen}\textcolor{white}{\textbf{ID}} & \cellcolor{sigagreen}\textcolor{white}{\textbf{Requisito}} & \cellcolor{sigagreen}\textcolor{white}{\textbf{MoSCoW}} \\ \hline
RF-01 & Monitorear el estado ambiental del aula & Must \\ \hline
RF-02 & Controlar remotamente los dispositivos del aula & Must \\ \hline
RF-03 & Detectar autom\'aticamente la ocupaci\'on del aula & Must \\ \hline
RF-04 & Controlar remotamente los proyectores & Must \\ \hline
RF-05 & Ajustar remotamente la climatizaci\'on & Must \\ \hline
RF-06 & Integrar el sistema con la videovigilancia existente & Should \\ \hline
RF-07 & Consultar un panel de control centralizado & Must \\ \hline
RF-08 & Generar y enviar alertas por anomal\'ias & Must \\ \hline
RF-09 & Predecir fallos de equipos mediante IA & Should \\ \hline
RF-10 & Registrar el historial de fallas y mantenimientos & Must \\ \hline
RF-11 & Notificar de forma inmediata las fallas cr\'iticas & Must \\ \hline
RF-12 & Gestionar las solicitudes de mantenimiento & Must \\ \hline
RF-13 & Controlar autom\'aticamente el apagado energ\'etico & Must \\ \hline
RF-14 & Analizar patrones de consumo energ\'etico & Should \\ \hline
RF-15 & Programar horarios automatizados de encendido y apagado & Must \\ \hline
RF-16 & Apagar autom\'aticamente equipos de aulas desocupadas & Must \\ \hline
RF-17 & Generar reportes estad\'isticos peri\'odicos & Should \\ \hline
RF-18 & Exportar reportes en PDF y Excel & Should \\ \hline
RF-19 & Gestionar acceso diferenciado por roles & Must \\ \hline
RF-20 & Registrar el historial de ocupaci\'on de aulas & Should \\ \hline
RF-21 & Notificar equipos encendidos fuera de horario & Must \\ \hline
RF-22 & Monitorear la conectividad de la red IoT & Must \\ \hline
RF-23 & Registrar bit\'acora de acciones de usuario & Must \\ \hline
RF-24 & Exportar los datos personales del usuario a solicitud & Should \\ \hline
RF-25 & Rectificar los datos personales del usuario & Should \\ \hline
\end{tabular}
\caption{Cat\'alogo de requisitos funcionales y su prioridad}
\label{tab:catalogo-rf}
\end{table}

A ellos se suman los seis requisitos funcionales de los componentes de inteligencia artificial (RF-IA-01 a RF-IA-06), que refinan RF-09 y RF-14 y se detallan en \S7.

\subsection{Requisitos no funcionales y cobertura del modelo de calidad}

Diecisiete requisitos no funcionales cuantificados (Tabla~\ref{tab:catalogo-rnf}), que cubren las nueve caracter\'isticas de ISO/IEC 25010:2023 \cite{iso25010_2023}. La columna de estado de validaci\'on es deliberadamente expl\'icita: distingue lo validado en campo de lo que sigue siendo un umbral propuesto sin confirmar, y esa distinci\'on se mantiene en la versi\'on final en lugar de suavizarse.

\begin{table}[H]
\centering
\small
\begin{tabular}{|l|p{5.0cm}|l|l|p{3.3cm}|}
\hline
\cellcolor{sigagreen}\textcolor{white}{\textbf{ID}} & \cellcolor{sigagreen}\textcolor{white}{\textbf{Requisito}} & \cellcolor{sigagreen}\textcolor{white}{\textbf{Caracter\'istica}} & \cellcolor{sigagreen}\textcolor{white}{\textbf{Prior.}} & \cellcolor{sigagreen}\textcolor{white}{\textbf{Validaci\'on}} \\ \hline
NFR-01 & Alertas entregadas en $\leq$60 s & Efic. de desempe\~no & Must & Validado en 1B \\ \hline
NFR-02 & Disponibilidad $\geq$99\,\% mensual & Fiabilidad & Must & Validado en 1B \\ \hline
NFR-03 & Cifrado AES-256 y acceso autenticado & Seguridad & Must & Validado en 1B \\ \hline
NFR-04 & Tarea b\'asica en $\leq$2 min sin capacitaci\'on & Interacci\'on de usuario & Should & Validado en 1B \\ \hline
NFR-05 & Operaci\'on en 3 navegadores, 360--1920 px & Compatibilidad & Must & Validado en 1B \\ \hline
NFR-06 & Cobertura de pruebas $\geq$70\,\% & Mantenibilidad & Should & Validado en 1B \\ \hline
NFR-07 & Acierto de ocupaci\'on $\geq$95\,\% & Adecuaci\'on funcional & Must & No verificado \\ \hline
NFR-08 & Sincronizaci\'on de horario sin p\'erdida & Compatibilidad & Must & Parcial (EV-10) \\ \hline
NFR-09 & Nuevo sensor en $\leq$8 h-persona & Flexibilidad & Could & No verificado \\ \hline
NFR-10 & Explicabilidad de la predicci\'on de RF-09 & Interacci\'on de usuario & Must & No verificado \\ \hline
NFR-11 & Derecho de acceso en $\leq$15 d\'ias h\'abiles & Seguridad / legal & Must & Normativo \\ \hline
NFR-12 & Alerta ante acceso no autorizado $\leq$5 min & Seguridad & Must & No verificado \\ \hline
NFR-13 & Reanudaci\'on tras ca\'ida de gateway $\leq$5 min & Fiabilidad & Should & No verificado \\ \hline
NFR-14 & Diagn\'ostico de falla $\leq$15 min & Mantenibilidad & Should & Parcial (EV-12) \\ \hline
NFR-15 & Operaci\'on con ancho de banda $\geq$1 Mbps & Flexibilidad & Should & No verificado \\ \hline
NFR-16 & Aviso expl\'icito de falta de conectividad & Fiabilidad & Must & Campo (EV-15) \\ \hline
RNF-17 & Apagado autom\'atico seguro & Seguridad f\'isica & Must & Borrador PE5 \\ \hline
\end{tabular}
\caption{Requisitos no funcionales, caracter\'istica de calidad, prioridad y estado de validaci\'on}
\label{tab:catalogo-rnf}
\end{table}

Nueve de los diecisiete siguen sin validaci\'on de campo. La causa est\'a documentada en el propio ERS: el guion de entrevista aplicado no inclu\'ia preguntas dirigidas a confirmar umbrales cuantitativos, de modo que ning\'un participante los discuti\'o. Declararlos como ``no verificado'' en lugar de darlos por validados es una decisi\'on deliberada, y es lo que permite que la m\'etrica de verificabilidad (\S8) se lea sin ambig\"uedad: un requisito puede ser \emph{verificable} (tener umbral, unidad y m\'etodo) sin estar todav\'ia \emph{verificado}.

\subsection{Actores y cobertura}

El sistema define nueve actores: cinco humanos y cuatro de sistema. La \'ultima columna de la Tabla~\ref{tab:actores} es la que sustenta la m\'etrica M1c (\S8) y, le\'ida junto a la anterior, expone la limitaci\'on de elicitaci\'on que el proyecto declara: hay actores que participan intensamente en el sistema sin haber sido nunca fuente de un requisito.

\begin{table}[H]
\centering
\small
\begin{tabular}{|p{4.6cm}|l|c|c|}
\hline
\cellcolor{sigagreen}\textcolor{white}{\textbf{Actor}} & \cellcolor{sigagreen}\textcolor{white}{\textbf{Tipo}} & \cellcolor{sigagreen}\textcolor{white}{\textbf{RF que origina}} & \cellcolor{sigagreen}\textcolor{white}{\textbf{CU en que participa}} \\ \hline
Personal de Infraestructura y Mantenimiento & Humano & 16 & 10 \\ \hline
Docente & Humano & 9 & 4 \\ \hline
Personal de TI & Humano & 0 & 6 \\ \hline
Personal Administrativo & Humano & 0 & 7 \\ \hline
Autoridades de la Facultad & Humano & 0 & 5 \\ \hline
Sensores / Gateway IoT & Sistema & 0 & 7 \\ \hline
Sistema de Horario Acad\'emico & Sistema & 0 & 3 \\ \hline
Servicio de Notificaciones & Sistema & 0 & 2 \\ \hline
Sistema de Videovigilancia & Sistema & 0 & 1 \\ \hline
\end{tabular}
\caption{Actores del sistema, requisitos que originan y casos de uso en que participan}
\label{tab:actores}
\end{table}

El diagrama general de casos de uso (Figura~\ref{fig:usecase-general}) es el artefacto sobre el que se cont\'o la \'ultima columna de la Tabla~\ref{tab:actores}: cada asociaci\'on actor--caso de uso del diagrama es una unidad de ese conteo, de modo que la m\'etrica M1c se puede reproducir sin m\'as insumo que la figura.

\begin{figure}[H]
\centering
\makebox[\textwidth][c]{\includegraphics[width=1.25\textwidth]{usecase_general.png}}
\caption{Diagrama general de casos de uso del sistema SIGA}
\label{fig:usecase-general}
\end{figure}

\subsection{Historias de usuario}

Los diecisiete requisitos de prioridad Must tienen historia en formato Connextra con un escenario de aceptaci\'on Gherkin. Las diecisiete se reescribieron por completo en la versi\'on 4.0 (defectos D-10 y D-13). La Tabla~\ref{tab:historias} resume el rol y el beneficio de cada una; el escenario completo est\'a en la \S3.6 del ERS.

\begin{table}[H]
\centering
\small
\begin{tabular}{|l|l|p{2.9cm}|p{6.2cm}|}
\hline
\cellcolor{sigagreen}\textcolor{white}{\textbf{HU}} & \cellcolor{sigagreen}\textcolor{white}{\textbf{RF}} & \cellcolor{sigagreen}\textcolor{white}{\textbf{Rol}} & \cellcolor{sigagreen}\textcolor{white}{\textbf{Beneficio buscado}} \\ \hline
HU-01 & RF-01 & Infraestructura & Conocer el estado de las aulas sin recorrerlas \\ \hline
HU-02 & RF-02 & Infraestructura & No desplazarse hasta el aula para operar equipos \\ \hline
HU-03 & RF-03 & Infraestructura & Activar el ahorro sin revisar c\'amaras manualmente \\ \hline
HU-04 & RF-04 & Infraestructura & Resolver la falla m\'as frecuente sin interrumpir la clase \\ \hline
HU-05 & RF-05 & Infraestructura & Mantener condiciones adecuadas sin desplazarse \\ \hline
HU-06 & RF-07 & Infraestructura & Detectar de un vistazo qu\'e aula necesita atenci\'on \\ \hline
HU-07 & RF-08 & Infraestructura & Enterarse de la anomal\'ia sin vigilar el panel \\ \hline
HU-08 & RF-10 & Infraestructura & Sustentar los informes al decanato con datos \\ \hline
HU-09 & RF-11 & Docente & Que la interrupci\'on se resuelva sin perder la hora \\ \hline
HU-10 & RF-12 & Docente & Saber si el reporte fue recibido y en qu\'e punto est\'a \\ \hline
HU-11 & RF-13 & Infraestructura & No depender de que alguien recorra el edificio \\ \hline
HU-12 & RF-15 & Infraestructura & Que el ahorro no dependa de una acci\'on diaria \\ \hline
HU-13 & RF-16 & Infraestructura & Eliminar el desperdicio energ\'etico nocturno \\ \hline
HU-14 & RF-19 & Personal de TI & Desplegar en producci\'on sin exponer operaciones cr\'iticas \\ \hline
HU-15 & RF-21 & Infraestructura & Corregir el desperdicio que motiv\'o el proyecto \\ \hline
HU-16 & RF-22 & Personal de TI & No decidir con datos desactualizados sin advertirlo \\ \hline
HU-17 & RF-23 & Personal de TI & Auditar qui\'en hizo qu\'e y cu\'ando \\ \hline
\end{tabular}
\caption{Historias de usuario de los requisitos Must}
\label{tab:historias}
\end{table}

Los ocho requisitos de prioridad Should no llevan historia ni criterio BDD, porque la gu\'ia los exige \'unicamente para los Must. Esa ausencia est\'a marcada como tal en la matriz de trazabilidad, con su raz\'on, y no como una cadena rota.

\subsection{Requisitos legales}

El sistema trata datos personales, por lo que el ERS mapea cada obligaci\'on relevante de la Ley Org\'anica de Protecci\'on de Datos Personales del Ecuador \cite{asamblea2021lopdp} al requisito que la operacionaliza: el Art. 13 (derecho de acceso) origina RF-24 y NFR-11; el Art. 14 (rectificaci\'on) origina RF-25; el deber de responsabilidad demostrada de los Art. 10 lit. k y 47 num. 2 se apoya en la bit\'acora de RF-23; y la notificaci\'on de vulneraciones de los Art. 43 y 46 se apoya en NFR-12. La \S3.4 del ERS argumenta adem\'as por qu\'e el procesamiento de imagen para determinar presencia agregada no constituye dato biom\'etrico en los t\'erminos del Art. 4, lo que sustenta la clasificaci\'on del proyecto en Categor\'ia B de riesgo \'etico.

\subsection{Restricciones de dise\~no}

Dieciocho restricciones acotan el espacio de soluci\'on en cinco ejes: tecnol\'ogico (RD-01 a RD-06: protocolos IoT, red institucional, base de datos aprobada, compatibilidad de equipos, l\'imite de la integraci\'on con videovigilancia y ejecuci\'on en navegador), econ\'omico (RD-07 y RD-08: presupuesto de hardware y despliegue por aulas piloto), normativo (RD-09 a RD-11: protecci\'on de datos, acceso por roles y bit\'acora), operativo (RD-12 y RD-13: dependencia del personal de Infraestructura y del horario acad\'emico digital), temporal (RD-14 y RD-15) y de alcance (RD-16 a RD-18, que fijan por escrito lo que el sistema no har\'a).

\subsection{Qu\'e cambi\'o en la versi\'on 4.0}

\begin{table}[H]
\centering
\small
\begin{tabular}{|p{4.3cm}|p{8.4cm}|l|}
\hline
\cellcolor{sigagreen}\textcolor{white}{\textbf{Cambio}} & \cellcolor{sigagreen}\textcolor{white}{\textbf{Motivo}} & \cellcolor{sigagreen}\textcolor{white}{\textbf{Defecto}} \\ \hline
RNF-17, seguridad f\'isica & El modelo de calidad no cubr\'ia \emph{Safety} pese a que el sistema emite comandos de apagado sobre equipos f\'isicos & D-04 \\ \hline
CU-17, derechos sobre datos & RF-24 y RF-25 colgaban de CU-11, cuyo prop\'osito es la gesti\'on de acceso, no el ejercicio de derechos & D-07 \\ \hline
\S4.11, DFD de nivel 1 & Ning\'un modelo describ\'ia el flujo de datos a nivel de proceso, eslab\'on exigido por la cadena de trazabilidad & D-08 \\ \hline
Prioridad de 13 RNF & Solo cuatro RNF ten\'ian prioridad MoSCoW declarada & D-02 \\ \hline
Reescritura de 17 historias & Connextra malformado y cl\'ausula \emph{Cuando} sin evento observable & D-10, D-13 \\ \hline
\S6.1, alcance del MVP & Las secciones 6.1 y 6.4 declaraban conjuntos distintos de RF; solo siete coincid\'ian & D-12 \\ \hline
Reclasificaci\'on de NFR-05 y NFR-15 & Usaban \emph{Portabilidad}, caracter\'istica que ISO/IEC 25010:2023 sustituy\'o por \emph{Flexibilidad} & D-14 \\ \hline
\end{tabular}
\caption{Cambios incorporados en la versi\'on 4.0 del ERS}
\label{tab:cambios-v4}
\end{table}

Los siete cambios provienen de la auditor\'ia de calidad (\S5), no de una revisi\'on editorial: cada uno cierra un defecto registrado con su ubicaci\'on exacta en el documento.
